\documentclass[11pt,a4paper]{article}

\usepackage[utf8]{inputenc}
\usepackage[T1]{fontenc}
\usepackage[french]{babel}
\usepackage{lmodern}
\usepackage[a4paper,margin=2.5cm]{geometry}
\usepackage{xcolor}
\usepackage{graphicx}
\usepackage{booktabs}
\usepackage{longtable}
\usepackage{array}
\usepackage{enumitem}
\usepackage{fancyhdr}
\usepackage{listings}
\usepackage{tikz}
\usetikzlibrary{positioning,arrows.meta,shapes.geometric}
\usepackage[colorlinks=true,linkcolor=blue!50!black,urlcolor=blue!60!black]{hyperref}

\definecolor{accent}{HTML}{1F4E79}
\definecolor{lightgray}{HTML}{F2F2F2}

\lstset{
  basicstyle=\ttfamily\footnotesize,
  backgroundcolor=\color{lightgray},
  frame=single,
  breaklines=true,
  columns=fullflexible,
  keepspaces=true
}

\pagestyle{fancy}
\fancyhf{}
\lhead{Cahier des charges -- Détection de la langue des signes}
\rhead{Deep Learning Avancé 2026--2027}
\cfoot{\thepage}
\renewcommand{\headrulewidth}{0.4pt}

\newcommand{\todo}[1]{\textcolor{red}{[#1]}}
\newcommand{\req}[1]{\textbf{#1}}

\begin{document}

%==================== PAGE DE GARDE ====================
\begin{titlepage}
\centering
\vspace*{1cm}
{\large Cycle Ingénieur en Intelligence Artificielle \& Data Science -- 3\textsuperscript{ème} année\par}
\vspace{0.5cm}
{\large Module : Deep Learning Avancé -- Année universitaire 2026--2027\par}
\vspace{2cm}
{\color{accent}\rule{\linewidth}{1.5pt}}\\[0.6cm]
{\Huge\bfseries Cahier des charges\par}
\vspace{0.4cm}
{\LARGE Détection de la langue des signes\\en temps réel\par}
\vspace{0.4cm}
{\large Industrialisation d'une solution de vision par ordinateur et MLOps bout-en-bout\par}
\vspace{0.3cm}
{\color{accent}\rule{\linewidth}{1.5pt}}\\[1.5cm]

{\large \textbf{Axe thématique retenu : Axe 3} -- Détection d'objets en temps réel (YOLO)\par}
\vspace{0.5cm}
{\large Architecture découplée FastAPI + React + MLOps (DVC, MLflow, Docker, CI/CD)\par}

\vspace{2cm}
\begin{tabular}{ll}
\textbf{Équipe :} & \todo{Nom Prénom 1}\\
 & \todo{Nom Prénom 2}\\
 & \todo{Nom Prénom 3}\\[0.4cm]
\textbf{Enseignant :} & Haythem Ghazouani\\[0.4cm]
\textbf{Version :} & 1.0 -- \today\\
\textbf{Dépôt Git :} & \todo{URL du dépôt GitHub}
\end{tabular}

\vfill
{\small Document soumis à validation de l'enseignant (Jalon 1)}
\end{titlepage}

\tableofcontents
\newpage

%==================== 1. CONTEXTE ====================
\section{Contexte et problématique}

Près de 70 millions de personnes sourdes dans le monde utilisent une langue des signes comme langue première. Or, la majorité des interlocuteurs entendants ne la maîtrisent pas, ce qui crée une barrière de communication quotidienne (administration, santé, éducation, commerce).

Les systèmes de vision par ordinateur permettent aujourd'hui de reconnaître des gestes manuels à partir d'une simple caméra. Cependant, un modèle entraîné dans un notebook ne constitue pas un produit : il doit être versionné, tracé, empaqueté, exposé via une API, consommé par une interface fluide et supervisé en production.

\paragraph{Problématique.} Comment concevoir, entraîner et industrialiser un système capable de \textbf{détecter et reconnaître en temps réel les lettres de la langue des signes (alphabet dactylologique)} à partir d'une image ou d'un flux webcam, avec une latence suffisamment faible pour une utilisation interactive, tout en respectant les standards professionnels du MLOps ?

\section{Objectifs du projet}

\subsection{Objectif général}
Livrer un produit logiciel complet, reproductible et prêt pour la production : un détecteur d'alphabet de la langue des signes en temps réel, accessible via une application web.

\subsection{Objectifs spécifiques}
\begin{enumerate}[leftmargin=*]
  \item \textbf{Deep Learning} : entraîner un détecteur YOLO (YOLOv8 / YOLO11) par fine-tuning, le comparer à une baseline, et l'évaluer avec des métriques rigoureuses (mAP, précision, rappel, F1, matrice de confusion).
  \item \textbf{MLOps} : versionner données et modèles avec DVC, tracer toutes les expériences et gérer le cycle de vie du modèle avec MLflow, exporter en ONNX.
  \item \textbf{Backend} : exposer le modèle via une API REST FastAPI asynchrone, typée (Pydantic) et testée (pytest).
  \item \textbf{Frontend} : développer un client React interactif (drag \& drop, galerie d'exemples, webcam temps réel, rendu des boîtes englobantes).
  \item \textbf{DevOps} : conteneuriser (multi-stage) et orchestrer par Docker Compose ; automatiser lint, tests et build via GitHub Actions.
\end{enumerate}

%==================== 2. PERIMETRE ====================
\section{Périmètre fonctionnel}

\subsection{Dans le périmètre}
\begin{itemize}[leftmargin=*]
  \item Détection et classification de \textbf{signes statiques} correspondant aux lettres de l'alphabet (A--Z, selon le jeu de données retenu).
  \item Inférence sur image unique (upload) et sur flux webcam (échantillonnage de trames).
  \item Affichage de la boîte englobante, de la lettre prédite et du score de confiance.
  \item Mesure et affichage de la latence (client et serveur).
  \item Pipeline MLOps complet et reproductible.
\end{itemize}

\subsection{Hors périmètre (extensions éventuelles)}
\begin{itemize}[leftmargin=*]
  \item Reconnaissance de signes dynamiques (mots, phrases) nécessitant un modèle temporel (LSTM, Transformer vidéo). Certaines lettres comportent un mouvement (par exemple J et Z en ASL) et ne peuvent pas être correctement capturées par une seule image : ce point est assumé comme limite connue.
  \item Traduction complète langue des signes $\leftrightarrow$ langue parlée.
  \item Déploiement cloud public (envisageable si l'équipe est élargie à 4 membres).
\end{itemize}

%==================== 3. DONNEES ====================
\section{Données}

\subsection{Jeu de données principal}
Le jeu retenu est \textbf{American Sign Language Letters} (Roboflow Public Datasets, auteur David Lee), jeu de détection d'objets où chaque lettre est annotée par une boîte englobante.

\begin{center}
\begin{tabular}{@{}ll@{}}
\toprule
\textbf{Caractéristique} & \textbf{Valeur} \\ \midrule
Tâche & Détection d'objets (boîtes englobantes) \\
Classes & 26 (A--Z) \\
Volume (export v1) & 1\,728 images (train 1\,512 / valid 144 / test 72) \\
Licence & Domaine public (selon la page du dataset) \\
Formats d'export & YOLOv8, YOLO11, COCO JSON, Pascal VOC, etc. \\
Source & \url{https://public.roboflow.com/object-detection/american-sign-language-letters} \\ \bottomrule
\end{tabular}
\end{center}

\subsection{Jeux complémentaires}
\begin{itemize}[leftmargin=*]
  \item \textbf{ASL Alphabet} (Kaggle, \texttt{grassknoted/asl-alphabet}) : environ 87\,000 images 200$\times$200 en 29 classes (A--Z, \textit{del}, \textit{nothing}, \textit{space}). Pas de boîtes : utilisable pour une baseline de classification (CNN / transfer learning) et pour la classe \og nothing \fg{} (absence de main).
  \item \textbf{ArSL21L} (Arabic Sign Language Letters, 14\,202 images, 32 classes, boîtes englobantes, licence CC BY 4.0) : variante ou extension pour la langue des signes arabe.
\end{itemize}

\subsection{Stratégie de qualité des données}
\begin{itemize}[leftmargin=*]
  \item \textbf{Aucune fuite de données} : le découpage train/valid/test est figé \emph{avant} toute augmentation ou normalisation ; l'augmentation ne s'applique qu'au jeu d'entraînement. Les images quasi-doublons issues d'un même sujet ne doivent pas être réparties entre train et test.
  \item Analyse exploratoire : distribution des classes, taille des boîtes, diversité des fonds et des éclairages.
  \item Augmentations (mosaic, HSV, rotations légères, flou, échelle). \emph{Attention} : le retournement horizontal est à évaluer avec précaution car il peut modifier la sémantique de certains signes (latéralité de la main).
  \item Test de robustesse sur des images propres à l'équipe (webcam, éclairages et fonds variés) pour mesurer l'écart de domaine.
\end{itemize}

%==================== 4. EXIGENCES ====================
\section{Exigences fonctionnelles}

\begin{longtable}{@{}p{1.3cm}p{9cm}p{2.2cm}@{}}
\toprule
\textbf{ID} & \textbf{Exigence} & \textbf{Priorité} \\ \midrule
\endhead
EF-01 & L'utilisateur peut déposer une image par glisser-déposer et obtenir la détection. & Obligatoire \\
EF-02 & L'utilisateur peut choisir une image dans une galerie d'exemples préchargés (inférence en 1 clic). & Obligatoire \\
EF-03 & L'utilisateur peut activer la webcam et obtenir une détection en continu. & Obligatoire \\
EF-04 & Les boîtes englobantes, labels et scores sont superposés dans un canvas, avec une couleur distincte par classe. & Obligatoire \\
EF-05 & Les probabilités des meilleures classes sont affichées en barres de progression animées. & Obligatoire \\
EF-06 & La latence d'inférence (ms) est affichée côté client et côté serveur. & Obligatoire \\
EF-07 & L'utilisateur peut régler le seuil de confiance et le seuil IoU (NMS). & Recommandée \\
EF-08 & Historique de la séquence de lettres détectées (construction d'un mot) avec stabilisation temporelle (vote majoritaire sur $N$ trames). & Recommandée \\
EF-09 & Téléchargement de l'image annotée. & Optionnelle \\
EF-10 & Mode sombre / clair, interface responsive. & Optionnelle \\ \bottomrule
\end{longtable}

\section{Exigences non fonctionnelles}

\begin{longtable}{@{}p{1.3cm}p{9cm}p{2.2cm}@{}}
\toprule
\textbf{ID} & \textbf{Exigence} & \textbf{Catégorie} \\ \midrule
\endhead
ENF-01 & Latence d'inférence serveur cible $< 100$\,ms par image sur CPU avec ONNX Runtime (valeur cible, à valider par mesure). & Performance \\
ENF-02 & Fluidité webcam d'au moins 10 images/s perçues côté utilisateur. & Performance \\
ENF-03 & Démarrage complet via une seule commande : \texttt{docker-compose up --build}. & Reproductibilité \\
ENF-04 & Aucune donnée ni poids de modèle volumineux versionné dans Git (DVC uniquement). & Gouvernance \\
ENF-05 & Couverture de tests pytest sur les endpoints et la fonction d'inférence. & Qualité \\
ENF-06 & Code conforme à PEP8 (flake8/black) et ESLint. & Qualité \\
ENF-07 & Les images traitées par la webcam ne sont ni stockées ni journalisées côté serveur (respect de la vie privée). & Éthique / RGPD \\
ENF-08 & L'API renvoie des codes HTTP sémantiques (400, 422, 503). & Robustesse \\ \bottomrule
\end{longtable}

%==================== 5. ARCHITECTURE ====================
\section{Architecture technique}

\subsection{Vue d'ensemble}

\begin{center}
\begin{tikzpicture}[
  node distance=1.1cm and 1.6cm,
  box/.style={rectangle,rounded corners,draw=accent,fill=accent!8,thick,minimum width=3.2cm,minimum height=1cm,align=center,font=\small},
  arr/.style={-{Stealth},thick}
]
\node[box] (user) {Utilisateur\\Navigateur / Webcam};
\node[box,right=of user] (front) {Frontend\\React + Tailwind};
\node[box,right=of front] (api) {API REST\\FastAPI + Pydantic};
\node[box,below=of api] (inf) {Moteur d'inférence\\ONNX Runtime (YOLO)};
\node[box,below=of front] (mlflow) {MLflow\\Tracking + Registry};
\node[box,below=of mlflow] (dvc) {DVC\\Données + checkpoints};

\draw[arr] (user) -- (front);
\draw[arr] (front) -- node[above,font=\scriptsize]{HTTP / WS} (api);
\draw[arr] (api) -- (inf);
\draw[arr] (mlflow) -- node[above,font=\scriptsize]{modèle validé} (inf);
\draw[arr] (dvc) -- node[left,font=\scriptsize]{versionnement} (mlflow);
\end{tikzpicture}
\end{center}

\subsection{Choix technologiques}

\begin{center}
\begin{tabular}{@{}p{4cm}p{9cm}@{}}
\toprule
\textbf{Composant} & \textbf{Technologie et justification} \\ \midrule
Modèle & YOLOv8/YOLO11 (Ultralytics), fine-tuning à partir de poids pré-entraînés COCO ; variante \textit{nano} ou \textit{small} pour le temps réel. \\
Baseline & Classifieur CNN (EfficientNet ou ResNet, transfer learning) sur recadrages de mains, comparé au détecteur. \\
Optimisation & Export ONNX, inférence ONNX Runtime, éventuellement FP16/INT8. \\
Backend & FastAPI, Pydantic, Uvicorn ; endpoints asynchrones. \\
Frontend & React (Vite), TailwindCSS ou Shadcn/UI, Axios/Fetch, canvas HTML5, \texttt{getUserMedia}. \\
Tracking & MLflow (paramètres, métriques, artefacts, Model Registry Staging/Production). \\
Données & DVC avec stockage distant (Google Drive, S3 ou local). \\
Conteneurs & Docker multi-stage, Docker Compose. \\
CI/CD & GitHub Actions : flake8, black, ESLint, pytest, build des images. \\ \bottomrule
\end{tabular}
\end{center}

\subsection{Spécification de l'API}

\begin{center}
\begin{tabular}{@{}llp{7.2cm}@{}}
\toprule
\textbf{Méthode} & \textbf{Route} & \textbf{Rôle} \\ \midrule
GET & \texttt{/health} & État du serveur, présence du GPU, version du modèle chargé. \\
GET & \texttt{/model/metadata} & Nom, architecture, métriques de validation, liste des classes. \\
POST & \texttt{/predict} & Reçoit une image (\texttt{UploadFile}), renvoie les détections en JSON (classe, score, boîte, latence). \\
POST & \texttt{/predict/visualize} & Renvoie l'image annotée avec boîtes, labels et scores. \\ \bottomrule
\end{tabular}
\end{center}

Exemple de réponse de \texttt{/predict} (schéma Pydantic) :
\begin{lstlisting}
{
  "detections": [
    {"label": "A", "confidence": 0.94,
     "bbox": {"x1": 112, "y1": 80, "x2": 310, "y2": 301}}
  ],
  "inference_time_ms": 38.2,
  "model_version": "yolo11n-asl-v3"
}
\end{lstlisting}

Gestion des erreurs : \texttt{400} si le fichier n'est pas une image valide, \texttt{422} pour une validation Pydantic échouée, \texttt{503} si le modèle n'est pas encore chargé.

%==================== 6. MLOPS ====================
\section{Cycle de vie MLOps}

\subsection{Versioning avec DVC}
Les dossiers \texttt{data/raw}, \texttt{data/processed} et les checkpoints sont suivis par DVC ; seuls les fichiers \texttt{.dvc} et \texttt{dvc.lock} sont commités. Un pipeline \texttt{dvc.yaml} reproductible enchaîne : préparation $\rightarrow$ entraînement $\rightarrow$ évaluation $\rightarrow$ export ONNX.

\subsection{Suivi d'expériences avec MLflow}
Chaque exécution journalise : hyperparamètres (taille du modèle, résolution, epochs, batch, optimiseur, augmentations), métriques par epoch (loss, mAP@0.5, mAP@[.5:.95], précision, rappel), artefacts (courbes PR, matrice de confusion, poids). Le meilleur modèle est enregistré dans le Model Registry et promu de \textit{Staging} à \textit{Production} selon un seuil de validation.

\subsection{Expériences prévues}
\begin{center}
\begin{tabular}{@{}llp{6cm}@{}}
\toprule
\textbf{Exp.} & \textbf{Modèle} & \textbf{Objectif} \\ \midrule
E0 & CNN transfer learning (classification) & Baseline de référence. \\
E1 & YOLO nano, 416 px & Premier détecteur, temps réel. \\
E2 & YOLO small, 640 px & Gain de précision vs coût de latence. \\
E3 & YOLO + augmentations avancées & Étude d'ablation sur la robustesse. \\
E4 & Export ONNX FP32 / FP16 / INT8 & Compromis précision / latence. \\ \bottomrule
\end{tabular}
\end{center}

\subsection{Conteneurisation et CI/CD}
\begin{itemize}[leftmargin=*]
  \item \texttt{backend/Dockerfile} et \texttt{frontend/Dockerfile} en multi-stage (image finale légère).
  \item \texttt{docker-compose.yml} orchestrant frontend, backend (et MLflow).
  \item \texttt{ci-tests.yml} : lint (flake8, black, ESLint) et pytest à chaque \textit{push} / \textit{pull request}.
  \item \texttt{docker-build.yml} : vérification du build des images.
\end{itemize}

%==================== 7. EVALUATION ====================
\section{Critères d'évaluation du modèle}

\begin{center}
\begin{tabular}{@{}lp{6.5cm}l@{}}
\toprule
\textbf{Indicateur} & \textbf{Description} & \textbf{Cible indicative} \\ \midrule
mAP@0.5 & Qualité globale de détection & $\geq 0{,}90$ \\
mAP@[.5:.95] & Précision de localisation & $\geq 0{,}65$ \\
F1 par classe & Détection des confusions entre lettres proches (M/N, U/V, etc.) & $\geq 0{,}85$ \\
Latence ONNX (CPU) & Temps serveur par image & $< 100$\,ms \\
Robustesse webcam & Précision sur images propres à l'équipe & à mesurer \\ \bottomrule
\end{tabular}
\end{center}

Les cibles sont indicatives et seront ajustées après l'analyse exploratoire et les premiers entraînements ; l'écart entre les scores sur le jeu de test et les performances réelles en webcam sera discuté de manière critique dans le rapport.

%==================== 8. ORGANISATION ====================
\section{Organisation de l'équipe}

\begin{center}
\begin{tabular}{@{}p{3.5cm}p{3cm}p{6.5cm}@{}}
\toprule
\textbf{Membre} & \textbf{Rôle} & \textbf{Responsabilités principales} \\ \midrule
\todo{Nom 1} & Lead Deep Learning & Exploration des données, entraînement, ablations, métriques, export ONNX. \\
\todo{Nom 2} & Lead MLOps / DevOps & DVC, MLflow, Docker Compose, GitHub Actions, profilage. \\
\todo{Nom 3} & Lead Full-Stack & API FastAPI, tests pytest, frontend React, UX webcam. \\ \bottomrule
\end{tabular}
\end{center}

Chaque membre reste responsable de l'ensemble des livrables et sera interrogé individuellement lors de la soutenance.

\paragraph{Bonnes pratiques Git.} Branche \texttt{main} protégée ; branches \texttt{feature/*} ; \textit{pull requests} relues par un autre membre ; messages de commit conventionnels (\texttt{feat:}, \texttt{fix:}, \texttt{docs:}, \texttt{ci:}).

%==================== 9. PLANNING ====================
\section{Planning et jalons}

\begin{center}
\begin{tabular}{@{}llp{7.2cm}@{}}
\toprule
\textbf{Jalon} & \textbf{Échéance} & \textbf{Livrables} \\ \midrule
J1 & Fin semaine 2 & Cahier des charges, choix du dataset, dépôt GitHub initialisé, DVC configuré. \\
J2 & Fin semaine 4 & Pipeline de données, baseline, fine-tuning YOLO, runs MLflow, métriques cibles validées. \\
J3 & Fin semaine 6 & Export ONNX, API FastAPI avec Swagger, tests pytest, Dockerfile backend. \\
J4 & Dernière séance & Frontend connecté, \texttt{docker-compose up}, CI/CD au vert, rapport PDF (LaTeX), soutenance technique avec démonstration en direct. \\ \bottomrule
\end{tabular}
\end{center}

\paragraph{Principe directeur.} L'intégration logicielle (API et interface) démarre dès le Jalon 2 afin d'éviter le \og \textit{notebook-only syndrome} \fg.

%==================== 10. RISQUES ====================
\section{Analyse des risques}

\begin{longtable}{@{}p{4.3cm}p{1.8cm}p{6.6cm}@{}}
\toprule
\textbf{Risque} & \textbf{Gravité} & \textbf{Mitigation} \\ \midrule
\endhead
Jeu de données petit et peu diversifié (sur-apprentissage, fonds uniformes) & Élevée & Augmentations, jeux complémentaires, collecte d'images propres à l'équipe, test de robustesse webcam. \\
Confusion entre lettres visuellement proches & Moyenne & Analyse de la matrice de confusion, ajout de données ciblées, stabilisation temporelle. \\
Signes dynamiques (ex. J, Z) non captés par une image & Moyenne & Limite documentée ; extension possible par suivi de trajectoire des points de la main. \\
Latence trop élevée sur CPU & Moyenne & Modèle nano, ONNX Runtime, quantification, réduction de la résolution. \\
Fuite de données & Élevée & Split figé avant toute augmentation, contrôle des doublons. \\
Poids ou données committés sur Git & Moyenne & \texttt{.gitignore} strict, DVC, vérification en CI. \\
Intégration tardive front/back & Élevée & API minimale dès J2, contrat OpenAPI partagé. \\
Biais (teints de peau, morphologies, éclairages) & Moyenne & Évaluation par sous-groupes lorsque possible, diversification des données de test. \\ \bottomrule
\end{longtable}

%==================== 11. LIVRABLES ====================
\section{Livrables et critères d'acceptation}

\subsection{Livrables}
\begin{enumerate}[leftmargin=*]
  \item Dépôt Git structuré (arborescence imposée en annexe A).
  \item Données et modèles versionnés avec DVC ; historique des runs MLflow.
  \item Modèle exporté au format ONNX.
  \item API FastAPI documentée (\url{http://localhost:8000/docs}) et tests pytest.
  \item Frontend React connecté.
  \item \texttt{docker-compose.yml} fonctionnel ; pipelines GitHub Actions.
  \item \texttt{README.md} d'installation et rapport technique (LaTeX, PDF).
\end{enumerate}

\subsection{Critères d'acceptation}
\begin{itemize}[leftmargin=*]
  \item \texttt{docker-compose up --build} lance l'application complète sans intervention manuelle.
  \item Inférence réussie depuis le navigateur sur upload, galerie et webcam, avec rendu des boîtes et latence affichée.
  \item Pipeline CI au vert (lint, tests, build).
  \item Aucun fichier lourd dans l'historique Git ; \texttt{dvc pull} restitue les données.
  \item Les résultats du rapport sont retrouvables dans MLflow.
\end{itemize}

\section{Considérations éthiques}
Le système traite des images de personnes. Les flux webcam ne sont pas conservés (ENF-07). Le projet est présenté comme une \emph{aide à la communication} et non comme un traducteur fiable : ses limites (alphabet statique, variabilité entre langues des signes, biais des données) sont explicitées dans le rapport.

%==================== ANNEXE ====================
\appendix
\section{Arborescence du dépôt}
\begin{lstlisting}
sign-language-detection/
|-- .github/workflows/
|   |-- ci-tests.yml
|   `-- docker-build.yml
|-- data/                  # versionne via DVC
|   |-- raw/
|   `-- processed/
|-- src/
|   |-- data/dataset.py
|   |-- models/{architecture,train,evaluate}.py
|   `-- utils/
|-- backend/
|   |-- Dockerfile
|   |-- requirements.txt
|   |-- app/{main,schemas,inference}.py
|   `-- tests/test_api.py
|-- frontend/
|   |-- Dockerfile
|   |-- package.json
|   `-- src/{components,services}/
|-- docker-compose.yml
|-- dvc.yaml
|-- README.md
`-- rapport_technique.pdf
\end{lstlisting}

\section{Références des jeux de données}
\begin{itemize}[leftmargin=*]
  \item David Lee, \textit{American Sign Language Letters Dataset}, Roboflow, 2020.
  \item Kaggle, \textit{ASL Alphabet} (\texttt{grassknoted/asl-alphabet}).
  \item Batnasan et al., \textit{ArSL21L: Arabic Sign Language Letter Dataset}, 2022.
\end{itemize}

\end{document}
